\documentclass[11pt]{article}
\usepackage[a4paper,margin=1in]{geometry}
\usepackage{amsmath,amssymb,mathtools,bm}
\usepackage{enumitem,booktabs,array}
\usepackage{tcolorbox}
\usepackage{hyperref}
\usepackage[protrusion=true,expansion=false]{microtype}
\hypersetup{colorlinks=true,linkcolor=blue,urlcolor=blue,citecolor=blue}
\setlist[itemize]{topsep=2pt,itemsep=2pt}
\setlist[enumerate]{topsep=2pt,itemsep=2pt}
\newtcolorbox{keybox}{colback=blue!3!white,colframe=blue!40!black,boxrule=0.4pt,arc=1mm}
\newtcolorbox{alertbox}{colback=red!3!white,colframe=red!40!black,boxrule=0.4pt,arc=1mm}
\newtcolorbox{answerbox}{colback=green!3!white,colframe=green!40!black,boxrule=0.4pt,arc=1mm}
\newtcolorbox{drillbox}{colback=yellow!5!white,colframe=orange!60!black,boxrule=0.4pt,arc=1mm}
\newcommand{\EV}{\mathbb{E}}
\newcommand{\Var}{\mathrm{Var}}
\newcommand{\Cov}{\mathrm{Cov}}
\newcommand{\blank}[1]{\underline{\hspace{#1}}}

\title{W15-01: Kamiya, Kang, Kim, Milidonis \& Stulz (2021, JFE) \\
\large ``Risk Management, Firm Reputation, and the Impact of Successful Cyberattacks on Target Firms'' --- Active Recall Workbook}
\author{Banking \& Risk Management --- Exam Prep}
\date{\today}
\begin{document}\maketitle

\begin{keybox}
\textbf{Paper in one sentence.} Using a hand-collected sample of disclosed successful cyberattacks on US public firms (2005--2017), KKKMS show that attacks involving loss of personal financial information cause a significant negative announcement return ($-1.1$\% CAR), large declines in sales growth and reputation (customer satisfaction), and lead target firms to increase risk management (cyber insurance, IT spending, board risk committees) --- evidence that cyber risk is material and affects real decisions.

\textbf{Placement.} Canonical cyber-risk-as-corporate-risk paper. The empirical complement to theories of reputational-risk management; pairs with Akey-Lewellen-Liskovich-Schiller (W15-02) on reputation channels.
\end{keybox}

\section*{Round 1: Complete Walk-Through}

\subsection*{1.1 Data}
Hand-collected sample of 480+ successful cyberattacks on US public firms 2005--2017. Sources: Privacy Rights Clearinghouse, SEC 8-K filings, news searches. Classify attacks by whether personal financial information (PFI) was lost.

\subsection*{1.2 Event-study spec}
\[
\text{CAR}_i(t_1,t_2)=\sum_{\tau=t_1}^{t_2}\!\big(R_{i\tau}-\EV[R_{i\tau}\mid \text{FF4}]\big).
\]
Windows around attack-disclosure date. Split by PFI-loss vs.\ not.

\subsection*{1.3 Headline results}
\begin{itemize}
\item Average $[-1,1]$ CAR: $\sim -0.8$\% for all attacks.
\item Attacks with PFI loss: $\sim -1.1$\% CAR; without PFI: insignificant.
\item Sales growth declines by $\sim 1$--$3$ pp over subsequent 2--3 years.
\item Reputation (ACSI customer satisfaction) falls.
\item Target firms increase cyber insurance, IT security spending, and board risk committees post-attack.
\item CFO turnover increases; CEO turnover not.
\end{itemize}

\subsection*{1.4 Reputation channel}
\begin{itemize}
\item Direct financial costs of breaches are modest relative to market-cap loss.
\item The gap is explained by reputational damage --- lost customer trust, brand value.
\item Product-market customers discipline targets: sales decline in B2C firms, less in B2B.
\end{itemize}

\subsection*{1.5 Risk-management response}
\begin{itemize}
\item Firms learn: post-attack investments in security/insurance rise.
\item Board-level risk oversight emerges (cyber-risk committees).
\item Consistent with optimal response to revealed exposure.
\end{itemize}

\subsection*{1.6 Caveats}
\begin{alertbox}
\begin{itemize}
\item Disclosure is endogenous; undisclosed attacks are missing.
\item PFI-loss attacks may cluster in certain industries.
\item Short panel for some risk-management outcomes.
\end{itemize}
\end{alertbox}

\section*{Round 2: Level 1 Fill-in}
\begin{enumerate}
\item Sample: \blank{1cm} US attacks, \blank{1cm}--\blank{1cm}.
\item CAR for PFI-loss attacks: $\sim$\blank{1cm}\%.
\item Sales-growth impact: $-$\blank{1cm} pp over 2--3 years.
\item Reputation measure: \blank{1.5cm} customer-satisfaction index.
\item Post-attack risk mgmt: cyber \blank{1.5cm}, IT spending, board \blank{1.5cm}.
\end{enumerate}

\section*{Round 3: Level 2 Prompts}
\begin{enumerate}
\item Discuss the event-study and the PFI split.
\item Explain the reputation channel and why CAR $>$ direct cost.
\item How do firms respond post-attack? What does this imply for optimal risk management?
\item Relate KKKMS to Akey-Lewellen-Liskovich-Schiller (2026).
\end{enumerate}

\section*{Round 4: Minimal Scaffolding}
From a blank page: (i) sample, (ii) event study, (iii) PFI heterogeneity, (iv) reputation/sales, (v) risk-mgmt response.

\section*{Round 5: Blank Page}
Essay: cyberattacks as a test of corporate risk management --- evidence from KKKMS.

\vspace{6pt}
\begin{drillbox}
\textbf{Speed Drill.} (1) Sample size. (2) CAR for PFI attacks. (3) Sales-growth impact. (4) Reputation measure. (5) Post-attack risk-mgmt actions.
\end{drillbox}

\section*{Quick Reference \& Answer Key}
\begin{answerbox}
\begin{itemize}
\item \textbf{Sample.} $\sim$480 US disclosed cyberattacks 2005--17.
\item \textbf{CAR.} $-1.1$\% for PFI-loss attacks.
\item \textbf{Sales.} Down 1--3 pp over 2--3 years.
\item \textbf{Response.} Insurance, IT, board committees up.
\end{itemize}
\end{answerbox}

\begin{keybox}
\textbf{30-second exam pitch.} Kamiya, Kang, Kim, Milidonis \& Stulz (2021) hand-collect $\sim$480 disclosed successful cyberattacks on US public firms, 2005--2017, to study the material impact of cyber risk. Attacks involving loss of personal financial information generate a significant $\sim -1.1$\% $[-1,1]$ CAR; attacks without PFI loss do not move stocks. The market-cap loss exceeds direct breach costs, revealing a large reputation channel: customer satisfaction (ACSI) falls, and sales growth declines by 1--3 pp over subsequent years, especially in B2C firms. Target firms respond by increasing cyber insurance, IT security spending, and board-level risk oversight; CFO turnover rises. The paper is the canonical empirical demonstration that cyber risk is a first-order material risk, affecting stock price, product-market outcomes, and real risk-management decisions --- a modern analogue to operational-risk management in financial firms.
\end{keybox}

\end{document}
